
% ============================================================================
\clearpage
\section*{v4.10.0-r22 Source-Only: The Collapse State Typed Across the Main Body: Reference Frame, Photon Point and Horizon}
\addcontentsline{toc}{section}{v4.10.0-r22 Source-Only: The Collapse State Typed Across the Main Body: Reference Frame, Photon Point and Horizon}
\label{sec:v4100-r22-collapse-state}
% ============================================================================

\paragraph{Purpose.}
The collapse state \(\zminus=0\) is the author's earliest idea and carries the central theorem of the main body
(Theorem~\ref{thm:unity}).  On 2026-11-16 the author asked for it to be re-examined with present knowledge.  v4.10.0-r21.9 read four
passages; this layer types all \(93\) occurrences of \(\zminus\) in the main body against Convention~F (v4.9.0-r22) and the
corrections of v4.9.0-r23, r24 and r29, in the way r23 and r24 typed the main statements.  It also corrects one row of r21.9.  A
reading, not blind; no computation and no prediction.  The main body is not rewritten.  The notes are
\texttt{notes/\allowbreak research/\allowbreak v4\_10\_0\_r22/} and
\texttt{notes/\allowbreak research/\allowbreak collapse\_state\_typing\_2026\_11\_16/}, with the list of occurrences.

\noindent\textbf{Source-only status.}
This v4.10.0-r22 layer is source-only.  It does not build a PDF, does not assign a new DOI,
and does not authorize a public release.  The v5.1.0 Version DOI
\href{https://doi.org/10.5281/zenodo.24237112}{10.5281/zenodo.24237112} and the Concept DOI
\href{https://doi.org/10.5281/zenodo.19210728}{10.5281/zenodo.19210728} are unchanged by it.

\subsection*{v4.10.0-r22 results}
\label{subsec:v4100-r22-results}

\paragraph{Remark v4.10.0-r22.1 (the typing).}
The occurrences fall into five groups.
\begin{center}
\small
\begin{tabular}{@{}>{\raggedright\arraybackslash}p{0.20\textwidth}>{\raggedright\arraybackslash}p{0.44\textwidth}>{\raggedright\arraybackslash}p{0.28\textwidth}@{}}
\toprule
group & passages & \(\zminus=0\) there is\\
\midrule
1.~zero of the offset & Theorems~\ref{thm:frame-offset} (proof), \ref{thm:symmetric-frame}, \ref{thm:dark-offset}, \ref{thm:spring}, \ref{thm:dark-energy} (absorbing state, \(\Lambda=0\)), \ref{thm:hawking-triolic}; Proposition~\ref{prop:shared-vacuum-planck}; the observer position, the GUE clock in the photon frame, \(k(c)=0\), \(\theta_H=0\), the readout chapters & R, the reference frame (r24.3, r29.1(a), r61.1(b), r64)\\
2.~the invariant object & Theorem~\ref{thm:unity} and the remark after it; ``the Sphaera is a photon''; light as the unique fixed point and as the compensator output of the universe & the role of \(P\), ascribed to the point \(\zminus=0\)\\
3.~the horizon & the definition in \S\ref{sec:bh}, Proposition~\ref{prop:horizon-cascade-singularity}, Theorem~\ref{thm:core}, the remark on the black hole as inverted Sphaera & \(\zminus\to\infty\) at \(\beta\to1\): the place of \(P\) (r24.3, r29.1(b))\\
4.~no point & Planck floor and tower \(|\zminus|_n\) (Theorem~\ref{thm:origin}), Pauli, \(J\colon\zminus\mapsto-\zminus\), the sign of \(\operatorname{Re}\zminus\), measurement as projection, the withdrawn RH row & scale, orientation (\(C\), r24.1), or withdrawn\\
5.~direction of the cascade & ``the cascade projects everything back toward \(\zminus=0\)'', Theorem~\ref{thm:core}, Theorem~\ref{thm:dark-energy} & open\\
\bottomrule
\end{tabular}
\end{center}
Group~1 holds about two thirds of the occurrences.

\paragraph{Theoremlet v4.10.0-r22.2 (the coordinate runs from R to \(P\); reading).}
In the main body \(\zminus=0\) is R, the pure state with zero offset, and \(\zminus\to\infty\) at \(\beta\to1\) is \(P\), the maximally mixed
state (r24.3, r29.1(b)).  The coordinate thus runs along the same line as the rapidity of v4.10.0-r21.8(d): there
\(\phi=\operatorname{sech}2\eta_R\) with R at \(\eta_R=0\) and \(P\) at \(\eta_R=\infty\); here \(|b|=\operatorname{sech}2\eta\) (r24.3).  Reading \(\zminus\) as a
rapidity is an interpretation: the main body uses \(\zminus\) on changing scales (\(\pm\bsym\) as a velocity, \(\beta/2\), \(\infty\) at \(\beta\to1\);
compare v4.8.0-r18).

\paragraph{Remark v4.10.0-r22.3 (the collapse state as a null state).}
The name ``collapse state'' sits at the end \(\zminus=0\), which in group~1 is R.  Group~2 ascribes to the same point the role of the
frame-invariant object, which belongs to \(P\).  Group~3 places \(P\), without naming it so, at the other end, \(\zminus\to\infty\).  The null
class of r21.8, the answer to the author's question whether at the lowest level a field lies in the null state or not, is therefore
the horizon end of the main-body coordinate, not its origin.  This agrees with v4.9.0-r29.1(b) and with the reading of the guarded core
as the shadow of a horizon (exploration of 2026-11-15, v4.10.0-r15.2).

\paragraph{Correction to v4.10.0-r21.9.}
Row~3 of the table in r21.9 read ``we can never fully reach \(\zminus=0\)'' (Theorem~\ref{thm:dark-offset}) as a statement of reachability
and placed it at \(P\).  In context the sentence says that the offset \(\bsym\) of the matter frame remains as a residual term, and
v4.8.0-r18 reads ``permanent'' as non-vanishing: the matter frame has a lasting offset and is never the pure state.  That is R, not
\(P\).  The tension with r24.3 recorded in r21.9 does not exist.  The other rows of r21.9 stand.

\paragraph{Remark v4.10.0-r22.4 (what stays open).}
Group~5 is not decided by the main body.  Whether ``back toward \(\zminus=0\)'' leads to the coordinate point R or to the class \(P\) depends
on the kind of step: zero offset points to R; an absorbing, non-invertible Markov cascade points to \(P\), since by k220 only
non-invertible steps (projection, mixing, limit) reach \(P\).

\paragraph{What this does not give.}
No rewriting of the main body; the typing is a reading at the listed passages.  No new number and no test.  The question which end
the cascade of group~5 approaches is open.

\subsection*{v4.10.0-r22 status in the three frames}
\label{subsec:v4100-r22-frames}

\begin{center}
\small
\begin{tabular}{@{}>{\raggedright\arraybackslash}p{0.30\textwidth}>{\raggedright\arraybackslash}p{0.62\textwidth}@{}}
\toprule
item & M / T / P (and R)\\
\midrule
group~1 & R, the frame of the main body; M at \(+\bsym\), \(T|_R\) at \(-\bsym\) (r24.2)\\
groups~2 and 3 & P: role at \(\zminus=0\), place at \(\zminus\to\infty\)\\
group~4 & frame-free (scale) or an orientation code\\
\bottomrule
\end{tabular}
\end{center}

\subsection*{v4.10.0-r22 register entries}
\label{subsec:v4100-r22-register}

\begin{center}
\small
\begin{tabular}{@{}>{\raggedright\arraybackslash}p{0.34\textwidth}>{\raggedright\arraybackslash}p{0.58\textwidth}@{}}
\toprule
statement & status after v4.10.0-r22\\
\midrule
occurrences of \(\zminus\) in the main body & \(93\), typed in five groups (r22.1)\\
coordinate \(\zminus\) & runs from R (\(0\)) to \(P\) (\(\infty\)); a rapidity under one reading (r22.2)\\
collapse state as null state & the horizon end \(\zminus\to\infty\), not the origin (r22.3)\\
row~3 of r21.9 & corrected: R, not \(P\); no tension with r24.3\\
direction of the cascade & open (r22.4)\\
\bottomrule
\end{tabular}
\end{center}

\subsection*{v4.10.0-r22 ledger and guard update}
\label{subsec:v4100-r22-ledger-guard}

\[
R_{\rm chk}^{v4100\text{-}r22}=(1,0,0)
\]
for (typing recorded; no computation; the direction of the cascade open).

The layer spends no guard:
\[
R_{\rm glob}^{(v4100\text{-}r22,+)}=(0,1,1,1,1),
\qquad
\Gamma_{\rm guard}^{v4100\text{-}r22}=(1,1,1,1,1),
\qquad
\beta_4=1/256\ \text{unspent}.
\]
The layer does not imply exact finite \(C5\) zero closure, exact all-readout terminality, \(C6/T2\), a global photon-clock/RH theorem,
or global all-readout closure.
